\documentclass{article}
\usepackage[margin=1in]{geometry}
\usepackage{hyperref}
\usepackage{parskip}

\title{B201 Computer Science Lab\\Portfolio Report}
\author{Hemanth Chandrashekar}
\date{June 2026}

\begin{document}
\maketitle
\section{Portfolio links}
The following links provide access to my portfolio website and its corresponding Github repository.
\begin{itemize}
    \item \textbf{Portfolio Website:}\href{https://hemanth-chandrashekar.github.io}{https://hemanth-chandrashekar.github.io}
    \item \textbf{Github Repository:}\href {https://github.com/Hemanth-Chandrashekar/hemanth-chandrashekar.github.io}{https://github.com/Hemanth-Chandrashekar/hemanth-chandrashekar.github.io}
    \item Stepstone (2026) \textit{Werkstudent Web Developer - berbach GmbH}. Available at: \href{https://stepstone.de/stellenangebote--WERKSTUDENT-WEB-DEVELOPER-M-W-D-Berlin-Kreuzberg-berbach-GmbH-Gestaltung-und-Beratung--12732537-inline.html}{stepstone.de}
\end{itemize}
\section{Portfolio Structure}
The portfolio website is structured  in such a way that the visitor goes through a logical journey. The first section the visitor sees is the hero section , which gives a quick overview of who i am. This is followed by the second section where the visitor goes through the about me section which gives a detailed information of who i am. The projects section showcases my skills and finally I've added the contact section if someone is interested to contact me. This order was chosen because as it mirrors the natural flow of how a job interview would be - the employer would want to know who i am first , then understand my background and skills,then they see the proof of my skills and finally if interested they'd contact me. 

The repository structure flows like this , index.html contains the HTML structure and content of the website , style.css contains the code for the visual styling elements of the website , CV/ folder has my professional CV written using the LaTeX software followed by the assets/ folder which contains the actual image files used in the portfolio and finally the README.md file which contains the overview of the project. 
\section{Design Decisions}
I chose to go with a dark theme because firstly it gives that professional developer vibe where the elements pop and it improves the readability, it reduces eye strain and i personally feel it would stand out against the white portfolio and it resembles standard tools like VS code and Terminal showing the similarities with the developer culture 

The space animation was chosen because i have a personal interest in space and astronomy and i feel space with moving stars represents the feeling of technology moving forward and the star animation demonstrates some skills of my javascript and it creates a cinematic impression which makes the portfolio stand out 

The monospace font was chosen because its clean and readable in dark backgrounds with smooth and easy edges creating a consistent techy aesthetic and it resembles the style used in code editors
\newpage
\section{Tools and Technologies}
\begin{itemize}
    \item \textbf{HTML}- HTML is used to build the basic structure of the website 
    \item \textbf{CSS}-CSS is used to add visual design  elements 
    \item \textbf{JavaScript}-JavaScript was used to create the cinematic space animation in the background of the portfolio 
    \item \textbf{LaTeX}-LaTeX is used to build my professional CV 
    \item \textbf{Git}-Git was used to track changes, save progress through commits, and manage version control of the project 
    \item \textbf{GitHub}-GitHub was used to host the repository online and store the project code remotely
    \item \textbf{GitHub Pages}-GitHub pages was used to host and deploy the portfolio website live on the internet for free
    \item \textbf{VS Code}- VS Code was used as the code editor to write and edit all the project files 
    \item \textbf{Linux Shell}- The linux shell was used to create,navigate and manage project files and run Git commands
\end{itemize}
\section{Strengths,Weaknesses and Future improvements}

The strengths are that the website is made from scratch using HTML and CSS , well documented GitHub repo ,dark techy theme stands out against plain white portfolios and lastly the cinematic space animation

The weaknesses of the portfolio are that there are limited projects, limited CSS animations, no real work experience and quite basic JavaScript

In the future i plan to improve my portfolio by adding more projects ,acquiring better skills of CSS and gaining real work experience 







\end{document}
